\begin{definition}[Predecessor in a Peano System]
\label{def:predecessor-in-peano-system}
Let $(P,S,1)$ be a Peano system, and let $x,u \in P$. The element $u$ is
called a \emph{predecessor} of $x$ exactly when
\[
S(u)=x.
\]
\end{definition}
\end{definitionbox}
\LeanFormalizes{def:predecessor-in-peano-system}{lra-lean}{LRA.VolumeII.PeanoSystems.TexSkeleton}{PredecessorInPeanoSystem}{pending}

\begin{remark*}[Standard quantified statement]
For a Peano system $(P,S,1)$ and elements $x,u\in P$,
\[
S(u)=x.
\]
\end{remark*}

\begin{remark*}[Negated quantified statement]
\[
\begin{aligned}
&\text{Let } (P,S,1) \text{ be a Peano system and let } x,u \in P.\\
&u \text{ is not a predecessor of } x
\Longleftrightarrow
S(u) \neq x.
\end{aligned}
\]
\end{remark*}

\begin{remark*}[Interpretation]
A predecessor of $x$ is an element whose successor is $x$. This language
isolates the inverse-looking question attached to the successor map without
assuming that successor is globally invertible.
\end{remark*}

\begin{dependencies}
\begin{itemize}
  \item \hyperref[def:peano-system]{Peano System}
\end{itemize}
\end{dependencies}
\end{topicbox}

\begin{topicbox}{Generation by Successor}
\begin{exposition}
The Peano axioms describe a counting system generated from a distinguished
first element by the successor operation. The next theorems make that
generation principle explicit: every element is either the first element or a
successor, every non-initial element has a unique predecessor, and no element
can be its own successor.
\end{exposition}

